\section{Methodology}
\label{sec:methodology}

Our framework treats segmentation as one stage in a two-stage decision pipeline.
The first stage produces a prediction and an uncertainty estimate. The second
stage, the deferral policy, consumes both and outputs a binary accept/defer
decision per pixel. The two stages are modular: any uncertainty method can be
paired with any deferral policy. This separation lets us isolate the effect of
each component and answer the question that matters for deployment: which
combination of uncertainty source and deferral rule produces the best decisions?

Figure~\ref{fig:pipeline} illustrates the full pipeline. We formalize it in
four parts: the problem setting (Section~\ref{sec:formulation}), uncertainty
estimation (Section~\ref{sec:uncertainty}), the deferral framework
(Section~\ref{sec:deferral}), and calibration (Section~\ref{sec:calibration}).

\begin{figure}[t]
\centering
\includegraphics[width=\textwidth]{pipeline.png}
\caption{Overview of the two-stage decision pipeline. Stage 1 produces a
segmentation prediction and uncertainty estimate. Stage 2 applies a deferral
policy to produce an accept/defer decision per pixel. The two stages are
modular: any uncertainty method can pair with any deferral policy.}
\label{fig:pipeline}
\end{figure}

\subsection{Problem formulation}
\label{sec:formulation}

Let $\mathbf{x} \in \mathbb{R}^{H \times W \times C}$ be an input image and
$\mathbf{y} \in \{0, 1\}^{H \times W}$ the corresponding binary segmentation
mask (vessel vs.\ background). A segmentation model $f_\theta$ produces a
probability map $\hat{\mathbf{p}} = f_\theta(\mathbf{x}) \in [0, 1]^{H \times
W}$, where $\hat{p}_{ij}$ is the predicted probability that pixel $(i, j)$ is
a vessel. The hard prediction is $\hat{y}_{ij} = \mathbb{1}[\hat{p}_{ij} >
0.5]$.

An uncertainty estimation procedure augments this pipeline with an uncertainty
map $\mathbf{u} \in \mathbb{R}_{\geq 0}^{H \times W}$. We use $u_{ij}$ to
denote the uncertainty assigned to pixel $(i, j)$. Higher values of $u_{ij}$
should, ideally, correspond to higher probability of error, i.e.,
$\hat{y}_{ij} \neq y_{ij}$.

A \emph{deferral policy} $\delta: \mathbb{R}_{\geq 0}^{H \times W} \times
[0,1]^{H \times W} \to \{0, 1\}^{H \times W}$ maps the uncertainty and
prediction to a binary decision map, where $\delta_{ij} = 1$ means ``accept
the prediction at pixel $(i,j)$'' and $\delta_{ij} = 0$ means ``defer this
pixel to a human reviewer.'' The goal is a policy that accepts correct
predictions and defers incorrect ones, subject to a coverage constraint: the
fraction of accepted pixels, $\text{coverage} = \frac{1}{HW}\sum_{ij}
\delta_{ij}$, should remain high enough that deferral is practical.

We evaluate deferral quality by the \emph{error reduction ratio}:
\begin{equation}
  \text{ERR} = \frac{e_{\text{before}} - e_{\text{after}}}{e_{\text{before}}}
  \label{eq:err}
\end{equation}
where $e_{\text{before}}$ is the error rate on all pixels and
$e_{\text{after}}$ is the error rate on accepted pixels only. A good deferral
policy achieves high ERR at high coverage.

\subsection{Uncertainty estimation}
\label{sec:uncertainty}

\subsubsection{MC Dropout}

MC Dropout \citep{gal2016dropout} keeps dropout active at test time and
performs $T$ stochastic forward passes through the network. For input
$\mathbf{x}$, let $\hat{p}_{ij}^{(t)} = f_{\theta_t}(\mathbf{x})_{ij}$ be
the predicted probability at pixel $(i,j)$ on pass $t$, where $\theta_t$
denotes the weights with dropout mask $t$ applied. The mean prediction is
\begin{equation}
  \bar{p}_{ij} = \frac{1}{T} \sum_{t=1}^{T} \hat{p}_{ij}^{(t)}.
  \label{eq:mc_mean}
\end{equation}

We quantify uncertainty using mutual information (MI), which decomposes total
uncertainty into aleatoric and epistemic components. For a binary prediction:
\begin{equation}
  \text{MI}_{ij} = H(\bar{p}_{ij}) - \frac{1}{T} \sum_{t=1}^{T} H(\hat{p}_{ij}^{(t)})
  \label{eq:mi}
\end{equation}
where $H(p) = -p \log p - (1-p) \log(1-p)$ is the binary entropy. The first
term is the predictive entropy (total uncertainty); the second is the expected
entropy across passes (aleatoric uncertainty). Their difference, MI, isolates
epistemic uncertainty: disagreement among the stochastic passes about the
correct class.

MI is preferable to raw variance for binary outputs because variance is bounded
by $0.25$ and is a symmetric, single-peaked function of $\bar{p}$, which
conflates high uncertainty with predictions near $p = 0.5$ regardless of
whether the model's passes agree. MI, by contrast, is zero when all passes
agree (even if they all predict $p = 0.5$) and large when passes disagree.

We use $T = 30$ passes with dropout probability $p_{\text{drop}} = 0.3$,
applied as spatial dropout (Dropout2d) after the decoder and before the
segmentation head. Inference is batched in chunks of 5 passes to manage GPU
memory.

\subsubsection{Test-Time Augmentation}

TTA applies a set of deterministic geometric transformations
$\{\tau_k\}_{k=1}^{K}$ to the input, performs a forward pass on each
transformed input, applies the inverse transformation $\tau_k^{-1}$ to the
output, and aggregates:
\begin{equation}
  \bar{p}_{ij}^{\text{TTA}} = \frac{1}{K} \sum_{k=1}^{K} \left[\tau_k^{-1}
  \circ f_\theta \circ \tau_k\right](\mathbf{x})_{ij}.
  \label{eq:tta_mean}
\end{equation}

We use $K = 6$ transformations: identity, horizontal flip, vertical flip,
90$^\circ$ rotation, 180$^\circ$ rotation, and 270$^\circ$ rotation. These
are exact (no interpolation artifacts) and cover the symmetry group of
the square.

Uncertainty is the pixel-wise variance across augmented predictions:
\begin{equation}
  u_{ij}^{\text{TTA}} = \frac{1}{K} \sum_{k=1}^{K} \left(\hat{p}_{ij}^{(k)} -
  \bar{p}_{ij}^{\text{TTA}}\right)^2
  \label{eq:tta_var}
\end{equation}
where $\hat{p}_{ij}^{(k)} = [\tau_k^{-1} \circ f_\theta \circ
\tau_k](\mathbf{x})_{ij}$.

We also compute the entropy of the mean prediction as an additional uncertainty
signal:
\begin{equation}
  u_{ij}^{\text{TTA-ent}} = H(\bar{p}_{ij}^{\text{TTA}}).
  \label{eq:tta_entropy}
\end{equation}

The two methods probe different sources of instability. MC Dropout samples
weight perturbations (epistemic uncertainty); TTA tests whether predictions
are stable under input transforms that should not change the answer. A
vessel is a vessel regardless of image orientation. Section~\ref{sec:discussion}
discusses why this distinction matters for deferral quality.

\subsection{Deferral framework}
\label{sec:deferral}

Given an uncertainty map $\mathbf{u}$ and a prediction map $\hat{\mathbf{p}}$,
we define three deferral policies.

\subsubsection{Global threshold deferral}

The simplest policy defers all pixels whose uncertainty exceeds a fixed
threshold $\tau$:
\begin{equation}
  \delta_{ij}^{\text{global}} = \mathbb{1}[u_{ij} \leq \tau].
  \label{eq:global}
\end{equation}

The threshold $\tau$ is selected on a validation set by sweeping over
percentiles of the uncertainty distribution and choosing the value that
maximizes a deferral-specific F1 score:
\begin{equation}
  F1_{\text{def}} = \frac{2 \cdot \text{prec}_{\text{def}} \cdot
  \text{rec}_{\text{def}}}{\text{prec}_{\text{def}} +
  \text{rec}_{\text{def}}}
  \label{eq:def_f1}
\end{equation}
where $\text{prec}_{\text{def}}$ is the fraction of deferred pixels that are
actual errors, and $\text{rec}_{\text{def}}$ is the fraction of all error
pixels that are deferred.

Global thresholding has the advantage of simplicity and a single
interpretable parameter. Its weakness is that a single $\tau$ must work
across images of varying difficulty: an easy image has low uncertainty
everywhere, so few pixels are deferred, while a hard image may have most
of its uncertainty mass above $\tau$, causing excessive deferral.

\subsubsection{Image-adaptive threshold deferral}

To address the variable-difficulty problem, image-adaptive thresholding
computes a separate threshold for each image based on its own uncertainty
distribution:
\begin{equation}
  \tau_n = Q_\alpha(\{u_{ij} : (i,j) \in \text{image } n\})
  \label{eq:adaptive_threshold}
\end{equation}
where $Q_\alpha$ is the $\alpha$-th percentile. The percentile $\alpha$
is chosen on a validation set; at test time, each image's threshold is
its own $\alpha$-th uncertainty percentile.

The deferral rule is then:
\begin{equation}
  \delta_{ij}^{\text{adaptive}} = \mathbb{1}[u_{ij} \leq \tau_n]
  \quad \text{for pixel } (i,j) \text{ in image } n.
  \label{eq:adaptive}
\end{equation}

This guarantees a fixed fraction $\alpha / 100$ of pixels are deferred
per image, which may be desirable in clinical workflows with a fixed
review budget. The trade-off is that easy images have pixels deferred
unnecessarily (the bottom $1-\alpha$ fraction of uncertainty in an
easy image may all be correct), while hard images may still not defer
enough.

We fit $\alpha$ using one of two criteria on the validation set:
\begin{itemize}[leftmargin=*]
  \item \emph{Max-F1}: choose $\alpha$ to maximize $F1_{\text{def}}$.
  \item \emph{Coverage-Dice}: choose the highest coverage (lowest deferral
  rate) such that Dice on accepted pixels exceeds a clinical threshold
  (0.82 in our experiments).
\end{itemize}

\subsubsection{Confidence-aware deferral}

Global and adaptive thresholding use only the uncertainty map $\mathbf{u}$.
But the prediction map $\hat{\mathbf{p}}$ contains additional information:
a pixel with $\hat{p}_{ij} = 0.51$ is near the decision boundary and likely
to flip under small perturbations, while $\hat{p}_{ij} = 0.95$ is far from
the boundary even if its uncertainty is moderate.

We define a prediction confidence score:
\begin{equation}
  c_{ij} = 2 \left| \hat{p}_{ij} - 0.5 \right| \in [0, 1]
  \label{eq:confidence}
\end{equation}
which is 0 when the prediction is maximally ambiguous ($\hat{p} = 0.5$) and
1 when the model is fully committed ($\hat{p} \in \{0, 1\}$).

The confidence-aware deferral score combines uncertainty and confidence:
\begin{equation}
  s_{ij} = u_{ij} \cdot (1 - c_{ij})
  \label{eq:conf_score}
\end{equation}

Pixels with high uncertainty \emph{and} low confidence receive the highest
scores and are deferred first. Pixels that are uncertain but confident
(the model predicts $p = 0.95$ but with high variance across passes)
receive lower scores and are retained.

The deferral rule is:
\begin{equation}
  \delta_{ij}^{\text{conf}} = \mathbb{1}[s_{ij} \leq \tau_s]
  \label{eq:conf_deferral}
\end{equation}
where $\tau_s$ is again chosen on a validation set.

A concrete example clarifies this. Consider two pixels, both with
$u_{ij} = 0.05$. Pixel A has $\hat{p} = 0.52$ (barely classified as
vessel), so $c_A = 0.04$ and $s_A = 0.05 \times 0.96 = 0.048$. Pixel B
has $\hat{p} = 0.92$ (clearly vessel), so $c_B = 0.84$ and
$s_B = 0.05 \times 0.16 = 0.008$. Global thresholding treats them
identically; confidence-aware deferral correctly prioritizes pixel A
for deferral.

\subsection{Calibration via temperature scaling}
\label{sec:calibration}

Temperature scaling \citep{guo2017calibration} learns a single parameter
$T > 0$ on a validation set to rescale the model's logits:
\begin{equation}
  \hat{p}_{ij}^{\text{cal}} = \sigma\left(\frac{\text{logit}(\hat{p}_{ij})}{T}\right)
  \label{eq:temp_scaling}
\end{equation}
where $\sigma$ is the sigmoid function and $\text{logit}(p) = \log(p / (1-p))$.

When $T > 1$, the sigmoid is flattened and predictions move toward 0.5
(the model becomes less confident). When $T < 1$, predictions sharpen.
$T$ is optimized by minimizing BCE loss on validation logits:
\begin{equation}
  T^* = \arg\min_T \; -\frac{1}{N_{\text{val}}} \sum_{ij}
  \left[ y_{ij} \log \hat{p}_{ij}^{\text{cal}} +
  (1 - y_{ij}) \log (1 - \hat{p}_{ij}^{\text{cal}}) \right]
  \label{eq:temp_opt}
\end{equation}
using L-BFGS, with $T$ constrained to $[0.05, +\infty)$.

In our experiments, the learned temperatures are $T = 1.19$ (MC Dropout) and
$T = 1.35$ (TTA), both $> 1$, confirming mild overconfidence. ECE changes
are negligible (Section~\ref{sec:results_calibration}), and the effect on
deferral quality is neutral to negative.

Why would calibration fail to help deferral? Temperature scaling is a
monotonic transformation of predicted probabilities: it rescales magnitudes
without changing rankings. If pixel A has higher uncertainty than pixel B
before calibration, it still does after. Deferral policies that threshold on
uncertainty ranking are therefore invariant to temperature scaling. The only
indirect path is through the confidence score $c_{ij}$ in the
confidence-aware policy, since softening predictions shifts which pixels sit
near the decision boundary. In practice, this second-order effect is small.
Section~\ref{sec:results_calibration} presents the evidence.
